\section{Les politiques conjoncturelles à l'œuvre depuis 1970}

This chapter follows the politiques conjoncturelles France has run since the end of
the Bretton-Woods system in 1971, in three periods: 1970--1982, when each
relance ran into inflation and the fall of the franc; 1983--1998, when a rigueur
monétaire under changes fixes led Europe to a single currency; and 2007--2026, when a
succession of crises raised the endettement. Its aim, in the course's words, is to
«~comprendre les mécaniques des politiques conjoncturelles, leurs impacts et leurs
conditions d'efficacité pour comprendre la situation actuelle~» --- the situation the
next chapter discusses.\\

The tools are those of the previous chapter: relance and rigueur, the effet
multiplicateur, the politique monétaire and the politique budgétaire, the taux de
change. What this chapter adds is the European system they work in, from the Serpent
monétaire européen to the BCE, and the budget rules, from the traité de Maastricht to
the reformed Pacte de Stabilité et de Croissance (PSC). The course opens it with Søren Kierkegaard: «~La vie doit être vécue
en regardant vers l'avenir, mais elle ne peut être comprise qu'en se retournant vers le
passé.~»

\subsection{Repérage historique des déséquilibres}

Three measures set the scene: croissance, inflation and the public accounts. The
croissance is the variation du PIB en volume in France: entre 1950 et 2024
$\Rightarrow$ 4 crises (recul du PIB en volume) = 1975, 1993, 2009, 2020, après les
chocs de 1973 et de 2008~; 4 pics de croissance = 76-78, 1988, 2000, 2021. Its trend breaks in 1974: before, the taux de croissance oscillates around 5 to 6~\% a year;
after, around 2~\%, and it never returns to the levels of the 1950s and 1960s, apart
from the 2021 rebound.

\begin{remark}[Formulation du cours]
  «~Entre 1950 et 2024 $\Rightarrow$ 4 crises = 1973, 1993, 2008, 2020~».
\end{remark}

Inflation, measured yearly since 1956, peaks in the \textbf{Grande inflation}
(\textit{Great Inflation}) of the 1970s and early 1980s, and then falls: «~À partir de
1983, baisse de l'inflation de 10~\% à 5~\% puis à 3-4~\% et stabilisation à 2~\% à
partir de 92 jusqu'en 2019.~» It was 0,5~\% in 2020 and 2~\% in 2024.\\

The public accounts: «~La France connaît un déficit public depuis 1974~». Until 1974
the dépenses publiques and the recettes publiques, in \% of PIB, move together; since
then spending has stayed above revenue, and the gap widens at each crisis (1975, 1993,
2009, 2020). In 2023 the dépenses publiques were about 57~\% of PIB. The déficit public is read
through the pair «~Stabilisateurs / déficit structurel~»: one part follows the
conjoncture, the other is the déficit structurel, defined below with the revisions of
the PSC.

\subsection{1970-1982 : l'échec de la relance keynésienne}

Every relance of the period met the same obstacle, the external constraint: inflation
higher than that of France's partners, a demand that went to imports, and réserves de
change spent defending the franc until a dévaluation became unavoidable. That is where
the failure lies. France then still controlled
capital movements, so the mobilité parfaite des capitaux of the model did not hold; and
in the Mundell-Fleming model, with changes fixes and mobilité parfaite des capitaux, it
is a monetary relance that is ineffective, while a fiscal relance is effective --- the
money the central bank must create to hold the parity adds to the relance (the case
of changes fixes is set out with the triangle d'incompatibilité below).

\begin{remark}[Formulation du cours]
  The period is headed: «~Pourquoi une politique de relance keynésienne est inefficace
  en situation de changes fixes, de mobilité parfaite des capitaux et d'inflation
  galopante~? Illustration: Les années 1970-1982 en France~».
\end{remark}

\subsubsection{Contexte international 1 : choc monétaire}

The first shock is monetary: la fin des \textbf{accords de Bretton-Woods}
(\textit{Bretton Woods agreements}) en 1971 et de la fixité des changes par rapport au
dollar, lui-même convertible en or. What ended in 1971 is the following system.

\begin{itemize}
  \item Les accords de Bretton-Woods adoptés en 1944 mettent en place un système
    monétaire international de \textbf{changes fixes} (\textit{fixed exchange rates})
    basés sur le dollar américain, et l'US\$ a une valeur fixe de conversion en or.
    Cela était permis par le fait que les US disposaient à l'époque des trois quarts
    des réserves d'or du monde.
  \item Les changes étaient alors \textbf{fixes mais ajustables} (\textit{adjustable
    peg})~: un pays pouvait dévaluer ou réévaluer sa monnaie, pourvu qu'un plan
    d'ajustement fût mis en œuvre et que les autres pays offrent leur accord. Le Fonds
    monétaire international (FMI) était chargé de la supervision de ce mécanisme.
  \item Le système a toutefois connu des difficultés à partir des années 1960 quand le
    stock de dollars détenu par les non-résidents est devenu supérieur au stock d'or
    américain (phénomène du dollar glut) $\Rightarrow$ les États-Unis n'ont plus pu
    assurer la conversion des dollars en or, et ont annoncé la suspension du système.
\end{itemize}

Conséquence~: \textbf{flottement généralisé des monnaies} (\textit{generalised floating})
à partir de 73 $\Rightarrow$ \textbf{changes flexibles} (\textit{flexible exchange
rates}) $\Rightarrow$ instabilité des changes $\Rightarrow$ les Européens de la
Communauté économique européenne (CEE) souhaitent conserver une stabilité relative de
leurs monnaies, pour maintenir le processus d'intégration économique depuis le traité de
Rome de 1957 $\Rightarrow$ coopération monétaire pour limiter les fluctuations des taux
de change entre les monnaies des pays de la CEE afin de stabiliser les échanges
commerciaux et protéger la politique agricole commune (prix agricoles fixes).\\

Hence the \textbf{Serpent monétaire européen} (\textit{European currency snake}),
1972--1979:
\begin{itemize}
  \item système à $\pm$~2,25~\% autour d'une \textbf{parité pivot} (\textit{central
    rate}) bilatérale pour limiter aussi les fluctuations entre les monnaies européennes
    face au dollar, sans mécanisme de solidarité financière~;
  \item si une monnaie atteignait la limite de sa marge, les banques centrales des deux
    pays concernés devaient intervenir pour ramener le taux vers la parité pivot~;
  \item les parités pivots n'étaient pas fixes~: elles pouvaient être réévaluées ou
    dévaluées en cas de déséquilibre persistant, collectivement, lors de réunions des
    ministres des Finances (exemple~: en 1976, plusieurs pays, dont l'Italie et la
    France, ont dévalué leur monnaie pour soutenir leur compétitivité).
\end{itemize}

\begin{example}[La marge du franc face au deutsche Mark]
  Si la parité pivot entre le franc français et le deutsche Mark (DM) était de
  3,35~FRF/DEM, le taux de change pouvait varier entre 3,27~FRF/DEM (plancher) et
  3,43~FRF/DEM (plafond). These are the two ends of the band: $3{,}35 \times (1 -
  0{,}0225) \approx 3{,}27$ and $3{,}35 \times (1 + 0{,}0225) \approx 3{,}43$.
\end{example}

\subsubsection{Contexte international 2 : chocs pétroliers 73 et 79}

\begin{itemize}
  \item Les \textbf{chocs pétroliers} (\textit{oil shocks}) sont un \textbf{choc d'offre}
    (\textit{supply shock}) $\Rightarrow$ dépendance de la France au pétrole importé du
    Moyen-Orient (99~\% en 1973 pour chauffage et électricité et 60~\% pour la
    production d'électricité en France contre 40~\% pour l'électricité en Allemagne).
  \item La croissance est très dépendante des importations de pétrole~: en 1973, 70~\% de
    l'énergie consommée en France était importée.
  \item Choc pétrolier 2~: la chute du Shah d'Iran et la révolution islamique
    $\Rightarrow$ arrêt des exportations de 79 à 81 (plus les perturbations de la guerre
    Iran/Irak en 1980) $\Rightarrow$ pénurie $\Rightarrow$ nouveau doublement du prix
    du baril (hausse durable des prix~: le prix du baril passe de moins de 4~\$ à 13~\$
    puis 35~\$ en 1979).
  \item Les conséquences économiques du deuxième choc furent aggravées par les
    fragilités de la croissance post-1973 et par une inflation élevée déjà ancrée dans
    les économies occidentales.
\end{itemize}
Réactions = nouvelles politiques énergétiques nationales
(renforcement des réserves, sources alternatives, économies d'énergie et plan nucléaire
en France).

\subsubsection{Chocs pétroliers $\Rightarrow$ stagflation}

Hausse du pétrole + dépendance énergétique + essence $\Rightarrow$ three chains of
effects:
\begin{itemize}
  \item hausse des coûts pour les entreprises $\Rightarrow$ licenciements~;
  \item hausse des prix $\Rightarrow$ baisse de la consommation~;
  \item hausse des salaires (indexés sur l'inflation) $\Rightarrow$ hausse des coûts et
    hausse de l'inflation $\Rightarrow$ cercle vicieux.
\end{itemize}

\begin{definition}[Stagflation]
  Conséquence = \textbf{stagflation} = stagnation + inflation.
\end{definition}

In France:
\begin{itemize}
  \item baisse de la croissance (2~\% en 1974 contre 6~\% en 1973)~;
  \item nouveauté~: un chômage persistant (2,5~\% en 1970~; 2,7~\% en 1973~; 6,1~\% en
    1979)~; de 1973 à 1978, l'emploi total a diminué d'environ 120\,000 personnes, alors
    que le nombre d'actifs s'est accru de 550\,000 (baby boomers)~;
  \item inflation galopante (plus de 10~\% dès 1974).
\end{itemize}
En Allemagne, stagflation plus faible~: un taux de chômage passant de 0,7~\% en 1970 à
4,7~\% en 1979, avec un pic à 4,8~\% en 1975~; une inflation de 6-7~\%~; une croissance
de 0,9~\% en 1974 contre 5,1~\% en 1973.

\subsubsection{Politiques conjoncturelles en France : alternance relance/rigueur}

The decade alternates between the two policies.

\begin{itemize}
  \item \textbf{1974~: plan de rigueur} --- contrer l'inflation en baissant les
    liquidités (encadrement du crédit, hausse de $r$, hausse des impôts sur les
    entreprises). Ces mesures restrictives réduisent la profitabilité des entreprises
    $\Rightarrow$ 1\textsuperscript{re} récession depuis la Seconde Guerre mondiale et
    hausse du chômage.
  \item \textbf{1975~: plan de relance keynésien de Chirac} --- relance de
    l'investissement et de la consommation (20 milliards de francs~: grands travaux,
    hausse des dépenses publiques). Relance de la croissance, mais pas de baisse du
    chômage, ni de l'inflation.
\end{itemize}

\begin{definition}[Spirale inflationniste (\textit{inflationary spiral})]
  Spirale inflationniste (= \textbf{inflation auto-entretenue} (\textit{self-sustaining
  inflation})) par 1) les salaires indexés sur les prix et 2) la chute du franc
  $\Rightarrow$ hausse du coût des importations (et donc aussi du pétrole =
  \textbf{importation incompressible} (\textit{incompressible imports})).
\end{definition}

\begin{itemize}
  \item \textbf{1976-79~:} plan d'\textbf{austérité} (\textit{austerity}) (Barre)~:
    «~la France vit au-dessus de ses moyens~».
    \begin{itemize}
      \item Gel des prix et des salaires pour réduire l'inflation, avec dégel entre 76 et
        78 et une fin progressive pour les salaires en 1982~;
      \item réduction du déficit public par une hausse des prélèvements obligatoires
        (de 29~\% à 39,4~\% du PIB entre 1975 et 1979), notamment des cotisations de
        Sécurité sociale, et par une augmentation de la fiscalité (hausse de la TVA, de l'IS et des impôts
        sur les ménages aisés)~;
      \item puis, à partir de 77, limitation des crédits~;
      \item en 78 et 79~: hausse des dépenses en faveur de l'emploi, mais hausse des
        cotisations sociales $\Rightarrow$ faiblesse des investissements $\Rightarrow$
        croissance surestimée $\Rightarrow$ hausse du déficit public, et toujours une
        inflation élevée et du chômage.
    \end{itemize}
  \item \textbf{1981~: plan de relance keynésien massif} --- 110 mesures de dépenses
    publiques d'un montant de 10 milliards de francs (1~\% du PIB)~: hausse de 10~\% du
    salaire minimum interprofessionnel de croissance (SMIC), retraite à 60 ans, passage
    aux 39 heures, embauche de 55\,000 fonctionnaires dès le mois de juin~; le minimum
    vieillesse et les allocations handicapés, familiales et logement augmentent de
    20~\% à 25~\%.
\end{itemize}

\subsubsection{La relance keynésienne de 81 bénéficie surtout aux importations et aux services}

La croissance connaît un léger rebond en 1982 (+2,5~\%), mais elle baisse dès 1983~:

\begin{center}
  \small
  \begin{tabular}{lcccccc}
    \toprule
    & 1979 & 1980 & 1981 & 1982 & 1983 & 1984 \\
    \midrule
    Évolution du PIB en volume & 3,6~\% & 1,6~\% & 1,1~\% & 2,5~\% & 1,3~\% & 1,5~\% \\
    \bottomrule
  \end{tabular}
\end{center}

\begin{itemize}
  \item Cette croissance est portée quasi exclusivement par la consommation des ménages
    et non l'investissement des entreprises $\Rightarrow$ elle profite
    principalement aux produits importés $\Rightarrow$ hausse du déficit commercial de
    la France, qui double presque (+87~\%) entre 1981 et 1982 (de 50 à 93 milliards de
    francs).
  \item Parallèlement, la politique sociale du nouveau gouvernement creuse massivement
    les déficits. Le dernier budget du gouvernement Barre (pour 1980) affichait un
    déficit quasi nul (0,4~\% du PIB)~; il passe à 2,4~\% du PIB en 1981 puis 2,8~\% en
    1982, un niveau égalant le record de 1975.
\end{itemize}
The «~services~» of the title are the tertiarisation described next.

\subsubsection{Crise ouvrière et industrielle en France dans les années 80}

\begin{itemize}
  \item Nationalisations~: en 1981-1982, le gouvernement Mauroy nationalise
    Saint-Gobain, mais c'est insuffisant pour redresser la compétitivité.
  \item Plans de restructuration~: le plan textile de 1982 prévoit des aides pour
    moderniser le secteur, mais les résultats sont limités~; dans l'automobile, l'État
    impose des quotas d'importation pour limiter les voitures japonaises, mais sans
    succès durable.
  \item Comparaison avec d'autres pays européens~:
    \begin{itemize}
      \item Allemagne~: maintient sa compétitivité grâce à des machines-outils
        performantes et une main-d'œuvre qualifiée~;
      \item Italie~: se spécialise dans le textile haut de gamme (mode, luxe) et résiste
        mieux~;
      \item Espagne/Portugal~: bénéficient de coûts salariaux très bas et attirent les
        délocalisations.
    \end{itemize}
  \item Bilan~: cette reprise en 1982 ne profite pas aux ouvriers qualifiés.
    \begin{itemize}
      \item Pas de soutien de l'emploi industriel, alors que les chocs pétroliers ont
        fragilisé les industries énergivores (sidérurgie, chimie)~;
      \item hausse des licenciements liés à la baisse de la compétitivité, à la
        modernisation (hausse de la productivité) et à la délocalisation~: entre 1982 et
        les années suivantes, l'emploi dans les secteurs de la production matérielle
        (industrie et agriculture) s'effondre~; selon l'Insee, 2,8 millions d'emplois
        industriels ont disparu depuis 1982, notamment dans des secteurs clés comme le
        textile et l'automobile, qui employaient massivement des ouvriers~;
      \item faible reconversion liée à la tertiarisation de l'économie~: la France se
        tourne vers les services (banque, tourisme, etc.), délaissant les industries
        lourdes et peu rentables comme le textile et l'automobile (emplois surtout créés
        dans les services, commerce et restauration) $\Rightarrow$ chômage de longue
        durée $\Rightarrow$ le taux de chômage passe de 8~\% en 1981 à 10,3~\% en 1983.
    \end{itemize}
\end{itemize}

\subsubsection{Multiples dévaluations dans les années 70 en France, entrées/sorties du Serpent monétaire et attaques spéculatives en 1981-82}

\begin{itemize}
  \item Janvier 1974~: la \textbf{politique monétaire accommodante}
    (\textit{accommodative monetary policy}) --- des taux «~bas~» pour relancer
    l'économie, mais des taux réels négatifs --- $\Rightarrow$ \textbf{fuite des
    capitaux} (\textit{capital flight}) $\Rightarrow$ sortie du Serpent $\Rightarrow$
    dévaluation du franc $\Rightarrow$ retour dans le Serpent.
\end{itemize}

\begin{definition}[Dévaluation]
  Dévaluation = renégociation de nouvelles parités par rapport au cours pivot pour
  éviter une sortie du système (Serpent ou Système monétaire européen ensuite).
\end{definition}

\begin{itemize}
  \item Sorties/entrées du Serpent~: le franc est dévalué quatre fois entre 1974 et
    1976, pour un total de près de 20~\%.
  \item Mars 1976~: la Banque de France (BdF) ne peut plus défendre le franc
    (effondrement des réserves) $\Rightarrow$ sortie définitive du Serpent, en raison
    des déséquilibres induits par la relance Chirac de septembre 1975
    $\Rightarrow$ chute du gouvernement.
\end{itemize}

\begin{remark}[Formulation du cours]
  The exit is dated and explained as «~sortie définitive du Serpent en raison des
  déséquilibres induits par la relance économique de 1976 de Chirac~».
\end{remark}

\begin{itemize}
  \item La France reste en dehors du Serpent jusqu'à la création du \textbf{Système
    monétaire européen} (\textit{European Monetary System}, SME) en 1979.
\end{itemize}

The SME brings a unit of account, the ECU (European Currency Unit); a support fund, the
\textbf{Fonds européen de coopération monétaire} (\textit{European Monetary
Cooperation Fund}, FECOM); and provision for \textbf{réalignements} (\textit{realignments}) of
the parités in a crisis. It limits fluctuation between European currencies only, not
against the dollar, within a more structured institutional frame. Both share the band of
$\pm$~2,25~\%.

\begin{center}
  \small
  \begin{tabularx}{\textwidth}{>{\raggedright\arraybackslash}p{2.9cm}
      >{\raggedright\arraybackslash}X >{\raggedright\arraybackslash}X}
    \toprule
    Critère & Serpent monétaire (1972-1979) & SME (1979-1999) \\
    \midrule
    Ancrage & Fluctuation limitée par rapport au dollar et entre les monnaies
      européennes. & Fluctuation limitée entre monnaies européennes. \\
    \addlinespace
    Bande de fluctuation & $\pm$~2,25~\%. & $\pm$~2,25~\%, étendue à $\pm$~6~\% pour
      certains pays~; $\pm$~15~\% à partir d'août 1993. \\
    \addlinespace
    Unité de compte & Aucune unité commune. & ECU~: panier de monnaies pour calculer les
      parités et les dettes. \\
    \addlinespace
    Mécanisme de soutien & Aucun fonds de solidarité. & FECOM~: aide financière
      temporaire aux pays en difficulté. \\
    \addlinespace
    Rôle de l'ECU & Non utilisé. & L'ECU sert de référence pour les parités, les dettes
      et les réserves. \\
    \addlinespace
    Flexibilité & --- & Possibilité de réalignements (ajustements de parité) en cas de
      crise. \\
    \addlinespace
    Coopération politique & Limitée, coordination informelle. & Cadre institutionnel
      plus structuré, avec des sommets européens réguliers. \\
    \bottomrule
  \end{tabularx}
\end{center}

\begin{itemize}
  \item Séries d'\textbf{attaques spéculatives} (\textit{speculative attacks}) en
    1981-82 à la suite du plan Mauroy $\Rightarrow$ anticipations de dévaluations
    $\Rightarrow$ elles provoquent la dévaluation réelle $\Rightarrow$ 2 dévaluations
    pour rester dans le SME.
\end{itemize}
Why? The sections that follow explain it.

\subsubsection{Mécanismes : échec des politiques de relance keynésienne en régime d'inflation élevée}

The table chains the causes. Two of them carry inflation from abroad and between prices
and wages: the \textbf{inflation importée} (\textit{imported inflation}) of the choc
pétrolier, and the \textbf{spirale prix-salaires} (\textit{wage-price spiral}) that the
\textbf{indexation des salaires} (\textit{wage indexation}) sets off. The déficit commercial
then drains the \textbf{réserves de change} (\textit{foreign exchange reserves}), the
foreign currencies the central bank sells to hold its parity.

\begin{center}
  \small
  \begin{tabularx}{\textwidth}{>{\raggedright\arraybackslash}p{3.2cm}
      >{\raggedright\arraybackslash}X >{\raggedright\arraybackslash}p{4.2cm}}
    \toprule
    Cause & Effet & Conséquence \\
    \midrule
    Relance budgétaire (dépenses publiques) & Hausse de la demande globale $\to$
      pression sur les prix~; hausse des importations (déficits jumeaux). & Inflation
      accrue. \\
    \addlinespace
    Choc pétrolier (1973) & Hausse des coûts de production $\to$ inflation importée. &
      Déficit commercial. \\
    \addlinespace
    Indexation des salaires & Spirale prix-salaires $\to$ inflation auto-entretenue. &
      Perte de compétitivité. \\
    \addlinespace
    Déficit commercial & Baisse des réserves de change $\to$ pression à la dévaluation.
      & Dévaluation du franc. \\
    \addlinespace
    Dévaluation & Rendre les exportations moins chères, mais renchérir les importations.
      & Nouvelle inflation importée. \\
    \addlinespace
    Cercle vicieux & Relance $\to$ inflation $\to$ dévaluation $\to$ inflation $\to$
      nouvelle relance\ldots & Stagflation et échec des politiques. \\
    \bottomrule
  \end{tabularx}
\end{center}

\subsubsection{Pourquoi un différentiel d'inflation provoque une dépréciation de la monnaie nationale~?}

Contexte~: l'inflation française dépassait celle de l'Allemagne (environ 13~\% en 1974
contre 7~\% en Allemagne)~--- a \textbf{différentiel d'inflation} (\textit{inflation
differential}). Two reasons make it depreciate the franc.\\

\textbf{Raison 1~: fuite des capitaux.}
\begin{itemize}
  \item Inflation ou anticipation d'inflation $\Rightarrow$ perte/anticipation de perte
    (dépréciation) de valeur de la monnaie $\Rightarrow$ recherche de placements dans
    des monnaies plus «~fortes~» (US\$, DM) $\Rightarrow$ sortie de capitaux (ventes de
    monnaie nationale contre des devises) $\Rightarrow$ dépréciation de la monnaie
    nationale.
  \item Plus l'inflation était élevée, plus le franc perdait de sa valeur sur les
    marchés des changes. Les investisseurs étrangers se détournaient du franc, aggravant
    la pression à la baisse.
\end{itemize}

\textbf{Raison 2~: déséquilibre de la balance commerciale.}
\begin{itemize}
  \item Inflation $\Rightarrow$ baisse de la compétitivité relative des exportations
    $\Rightarrow$ baisse des exportations.
  \item Parallèlement, inflation importée (énergie) $\Rightarrow$ hausse du coût des
    importations.
  \item Si au total déficit commercial $\Rightarrow$ entrées de devises $<$ sorties de
    devises $\Rightarrow$ hausse de la demande de devises et baisse de la demande de
    monnaie nationale $\Rightarrow$ dépréciation de la monnaie nationale.
\end{itemize}

\subsubsection{Pourquoi les dévaluations entrainent des attaques spéculatives~?}

\begin{itemize}
  \item En régime de changes fixes, pour maintenir la parité dans le Serpent monétaire
    ou dans le SME, la BdF doit vendre des devises pour acheter des francs (pour
    réapprécier le franc) $\Rightarrow$ diminution des réserves de change $\Rightarrow$
    quand les réserves de change sont épuisées $\Rightarrow$ dévaluation obligatoire.
  \item Des dévaluations en réaction à des déséquilibres réels (déficit commercial,
    inflation élevée, perte de compétitivité industrielle).
  \item Perte de confiance dans la politique économique française, qui ne propose pas
    de solution crédible $\Rightarrow$ crise de confiance et attaques spéculatives. The
    attaque spéculative comes \emph{before} the dévaluation: that is the mécanisme de l'attaque
    spéculative.
  \item Mais les dévaluations augmentent aussi le coût des importations pour les
    entreprises et alimentent les anticipations de «~dévaluation en cascade~»
    $\Rightarrow$ attaques spéculatives multiples.
  \item Si un pays dévaluait sa monnaie pour booster ses exportations, cela pouvait
    déclencher une \textbf{course à la dévaluation} (\textit{competitive devaluation})
    entre pays, nuisant à la stabilité économique européenne $\Rightarrow$ nombreux
    réajustements des parités centrales (ex.~: réévaluations du mark et du florin
    néerlandais~; dévaluations de la lire et du franc français en 1981).
  \item Bilan~: instabilité financière européenne.
\end{itemize}

\begin{example}[Les attaques contre le franc en 1981]
  Malgré une hausse des taux d'intérêt (le taux à 3 mois a atteint 20~\% et le taux à
  10 ans est passé de 13,4~\% fin 1980 à 17,5~\% en juin 1981), entre février et mai
  1981, après les annonces des 110 mesures du plan Mauroy $\Rightarrow$ attaques
  spéculatives $\Rightarrow$ la BdF a dû dépenser 5 milliards de dollars pour soutenir
  le franc sur les marchés des changes.
\end{example}

What exposes a currency to an attaque spéculative is of three kinds.
\begin{itemize}
  \item \textbf{Déséquilibres économiques sous-jacents.}
    \begin{itemize}
      \item La balance des paiements est déficitaire~: le pays importe plus qu'il
        n'exporte, ce qui épuise ses réserves de change.
      \item L'inflation est élevée~: une monnaie perd de sa valeur si les prix
        intérieurs augmentent trop vite, réduisant son pouvoir d'achat.
      \item La dette publique est insoutenable~: si un État emprunte massivement en
        devises étrangères (comme en dollars), une dépréciation de sa monnaie augmente
        mécaniquement le coût de sa dette.
      \item Les taux d'intérêt sont trop bas~: une politique monétaire trop accommodante
        peut inciter à des sorties de capitaux vers des placements plus rémunérateurs
        ailleurs.
    \end{itemize}
  \item \textbf{Surévaluation de la monnaie} (\textit{overvaluation}).
    \begin{itemize}
      \item Une monnaie est trop forte si son taux de change ne reflète pas sa réalité
        économique (ex.~: un pays avec une croissance faible mais une monnaie stable
        grâce à des entrées de capitaux spéculatifs).
      \item Les spéculateurs anticipent alors une correction (dépréciation) et parient
        contre la monnaie en la vendant à découvert.
    \end{itemize}
  \item \textbf{Manque de crédibilité des autorités monétaires.}
    \begin{itemize}
      \item Politique monétaire incohérente~: si une banque centrale alterne entre
        relance et austérité, les marchés perdent confiance.
      \item Réserves de change insuffisantes~: si un pays n'a pas assez de devises pour
        racheter sa propre monnaie sur les marchés, il ne peut pas résister à une
        attaque.
      \item Instabilité politique~: coups d'État, élections contestées ou corruption
        sapent la confiance dans la monnaie. Exemple~: en 2019, l'Argentine a subi une
        attaque spéculative après des élections perçues comme défavorables aux réformes
        économiques.
    \end{itemize}
\end{itemize}

Countries protect themselves from attaques spéculatives in five ways:
\begin{itemize}
  \item accumuler des réserves de change (ex.~: la Chine détient plus de 3\,000
    milliards de dollars)~;
  \item éviter une surévaluation (ne pas maintenir artificiellement un taux de change
    trop fort)~;
  \item \textbf{contrôles des capitaux} (\textit{capital controls})~: limiter les
    sorties de devises, comme en Malaisie en 1998~; taxe Tobin~;
  \item politique monétaire crédible (taux d'intérêt élevés pour attirer les capitaux,
    mais risqué pour la croissance)~;
  \item réformes structurelles (réduire la dette, améliorer la compétitivité, lutter
    contre la corruption).
\end{itemize}

\subsubsection{Pourquoi la hausse des taux nominaux n'a pas limité l'inflation~?}

Five reasons, two of them concepts in their own right: the \textbf{crédibilité}
(\textit{credibility}) a policy lacked, which recurs through the chapter, and the
\textbf{rigidités structurelles} (\textit{structural rigidities}) of the economy.

\begin{center}
  \small
  \begin{tabularx}{\textwidth}{>{\raggedright\arraybackslash}p{3.6cm}
      >{\raggedright\arraybackslash}X}
    \toprule
    Raison & Explication \\
    \midrule
    Chocs externes & Chocs pétroliers et hausse des matières premières ont alimenté
      l'inflation par les coûts. \\
    \addlinespace
    Taux réels négatifs & L'inflation élevée a annulé l'effet restrictif des taux
      nominaux $\Rightarrow$ hausse des emprunts et baisse de l'épargne. \\
    \addlinespace
    Indexation des salaires & Les salaires suivaient l'inflation passée, créant un
      cercle vicieux $\Rightarrow$ hausse de la demande alors que l'offre est limitée
      par les coûts~; + anticipation d'inflation. \\
    \addlinespace
    Manque de crédibilité & Les politiques monétaires étaient incohérentes avec les
      politiques budgétaires et peu crédibles (pas de cible d'inflation pour ancrer les
      anticipations), limitant leur impact. \\
    \addlinespace
    Rigidités structurelles & Marché du travail et secteur public peu flexibles,
      retardant l'ajustement des prix~; entreprises publiques avec des prix fixes. \\
    \bottomrule
  \end{tabularx}
\end{center}

\subsubsection{Zoom théorique : triangle d'incompatibilité de Mundell-Fleming}

\begin{definition}[Triangle d'incompatibilité de Mundell-Fleming (\textit{impossible trinity})]
  Principe développé par Robert Mundell et Marcus Fleming d'après lequel il n'est pas
  possible pour un État de maintenir à la fois un taux de change fixe, la liberté de
  circulation des capitaux, et l'indépendance de la politique monétaire de la banque
  centrale.
\end{definition}

The three objectives are the vertices of the triangle: the régime de change fixe, the
indépendance de la politique monétaire, and the \textbf{mobilité parfaite des capitaux}
(\textit{perfect capital mobility}), the triangle's name for la liberté de circulation
des capitaux. «~Ce triangle veut dire que si un pays veut une politique monétaire
efficace (ex: baisse de $r$ pour relancer l'économie), il faut enlever un des deux
autres points~» --- «~sauf à avoir des réserves de change très élevées comme
l'Allemagne ou la Chine~!~»

\begin{figure}[h!]
  \centering
  \begin{tikzpicture}
    \coordinate (F) at (90:2.8);
    \coordinate (I) at (210:2.8);
    \coordinate (M) at (330:2.8);
    \draw[thick] (F) -- (I) -- (M) -- cycle;
    \node[above=3pt] at (F) {Régime de change fixe};
    \node[below left=2pt and -12pt, align=center] at (I) {Indépendance de la\\politique monétaire};
    \node[below right=2pt and -12pt, align=center] at (M) {Mobilité parfaite\\des capitaux};
    \node[left=4pt] at ($(F)!0.5!(I)$) {a};
    \node[right=4pt] at ($(F)!0.5!(M)$) {b};
    \node[below=4pt] at ($(I)!0.5!(M)$) {c};
  \end{tikzpicture}
  \caption*{\small Triangle d'incompatibilité de Mundell-Fleming. Each vertex is one of
    the three objectives, and each side a, b, c joins two of them: a country can hold
    the two ends of one side, never all three vertices at once.}
\end{figure}

\begin{example}[Les réserves de dollars de la Chine]
  Entre 2000 et 2014, la Chine a accumulé des réserves de change record (plus de
  4\,000 milliards de dollars en 2014). Cette accumulation était due à un excédent
  commercial chronique (exportations $>$ importations). Grâce à ces réserves, la Chine
  a pu~:
  \begin{itemize}
    \item financer ses importations de matières premières (pétrole, minerais) sans
      dépendre des marchés financiers internationaux~;
    \item stabiliser le yuan face au dollar, évitant une appréciation trop forte qui
      aurait nui à ses exportations~;
    \item investir à l'étranger (rachats d'entreprises, prêts aux pays en
      développement)~;
    \item résister aux chocs externes (comme la crise financière de 2008).
  \end{itemize}
\end{example}

\paragraph{Pourquoi en changes flexibles la politique budgétaire est moins efficace qu'en changes fixes~?}
The triangle also decides how far a fiscal relance goes.
\begin{itemize}
  \item En changes fixes~: si le gouvernement augmente ses dépenses, la demande globale
    monte, ce qui peut faire pression à la hausse sur les prix et les taux d'intérêt.
    Mais la banque centrale doit maintenir le taux de change fixe en ajustant l'offre de
    monnaie (achat/vente de devises). Conséquence~: si la banque centrale augmente
    l'offre de monnaie pour défendre le taux de change, cela stimule davantage
    l'économie (effet expansionniste supplémentaire). L'inflation peut augmenter, mais
    l'impact sur la croissance est positif et durable.
  \item En changes flexibles~: si le gouvernement augmente ses dépenses, la demande
    globale monte, ce qui peut faire monter les taux d'intérêt (si la banque centrale
    ne réagit pas), attirant des capitaux étrangers, et provoquer une appréciation du
    taux de change (car la demande de monnaie nationale augmente) $\Rightarrow$ les
    exportations deviennent plus chères (perte de compétitivité), les importations
    deviennent moins chères (les ménages et entreprises achètent plus à l'étranger)
    $\Rightarrow$ la demande globale se déplace des biens nationaux vers les biens
    étrangers $\to$ baisse de la production intérieure. L'effet expansionniste initial
    est partiellement annulé~: une partie de l'impact expansionniste est neutralisée
    par l'appréciation du taux de change $\Rightarrow$ effet multiplicateur moins fort
    en changes flexibles.
\end{itemize}
This is the result the subsection opened with. Under changes fixes a fiscal relance is
effective in the model; what fails is a monetary relance, as in January 1974, when low
rates brought fuite des capitaux and an exit from the Serpent. The relances of 1975 and
1981 failed through inflation, imports and the loss of réserves de change.

\subsubsection{Est-ce que l'Allemagne est confrontée au TMF~?}

Germany, inside the Serpent, kept an autonomous politique monétaire. How it escaped the
TMF (the triangle de Mundell-Fleming):
\begin{itemize}
  \item Culture anti-inflationniste $\Rightarrow$ politique monétaire autonome et
    rigoureuse~:
    \begin{itemize}
      \item traumatisme de l'hyperinflation des années 20 (60~\% en 1921~; 5\,200~\% en
        1922~; 16\,580\,000~\% en 1923) $\Rightarrow$ rigueur~;
      \item la Bundesbank était indépendante et priorisait la stabilité des prix au
        détriment de la croissance ou de l'emploi $\Rightarrow$ contrôle de l'inflation
        avec des cibles précises et des taux réels positifs (5-6~\%) (futur modèle de la
        BCE~!)~: environ 13~\% d'inflation en France en 1974 contre 7~\% en Allemagne.
    \end{itemize}
  \item Politique budgétaire restrictive~: l'Allemagne évite les déficits publics
    excessifs.
  \item Conséquence $\Rightarrow$ appréciation continue du DM $\Rightarrow$ confiance
    dans la monnaie~:
    \begin{itemize}
      \item après les chocs pétroliers, hausse des taux d'intérêt $\Rightarrow$ entrée
        de capitaux étrangers $\Rightarrow$ pression à la hausse sur le DM
        $\Rightarrow$ menace sur la parité fixe avec les autres monnaies européennes~;
      \item effet des exportations $\Rightarrow$ sources de devises à convertir
        $\Rightarrow$ hausse de la demande de DM $\Rightarrow$ appréciation~;
      \item et l'Allemagne, avec une inflation plus faible que ses partenaires (ex.~:
        France, Italie), voyait sa monnaie s'apprécier mécaniquement (car les autres
        pays dévaluaient pour compenser leur inflation).
    \end{itemize}
  \item Pour rester dans le Serpent $\Rightarrow$ rupture du triangle de Mundell-Fleming
    par l'ajustement marginal de son taux de change~:
    \begin{itemize}
      \item l'excédent commercial permet de soutenir le mark sans craindre une pénurie
        de devises (en 1974, l'Allemagne affichait un excédent de 10 milliards de
        dollars, tandis que la France avait un déficit de 6 milliards)~;
      \item appréciation du DM $\Rightarrow$ suppression du risque d'inflation importée~;
      \item $\Rightarrow$ achat de devises étrangères (des exportations) contre DM pour
        augmenter la quantité de DM $\Rightarrow$ limite l'appréciation~;
      \item $\Rightarrow$ rachat de monnaies européennes contre DM.
    \end{itemize}
  \item Bilan~: pas d'attaques spéculatives~; flexibilité relative dans le TMF (moins de
    contrainte car monnaie forte) $\Rightarrow$ devient le modèle de référence pour la
    stabilité économique et financière en Europe.
\end{itemize}

The Bundesbank, première banque centrale nationale (BCN) indépendante~:
\begin{itemize}
  \item créée en 1957, elle a renforcé son indépendance dans les années 70 vis-à-vis du
    gouvernement, une caractéristique qui deviendra emblématique de sa politique~;
  \item ses objectifs principaux étaient~:
    \begin{itemize}
      \item stabilité des prix~: priorité absolue, même au détriment de la croissance ou
        de l'emploi~;
      \item stabilité du deutsche Mark~: maintenir la force du DM face aux autres
        devises, notamment le dollar~;
      \item crédibilité internationale~: la Bundesbank est devenue un modèle de banque
        centrale rigoureuse, inspirant plus tard la BCE~;
      \item politique de «~\textbf{monétarisme pragmatique}~» (\textit{pragmatic
        monetarism})~: la Bundesbank a adapté les cibles monétaires en fonction des
        chocs économiques, évitant une approche dogmatique.
    \end{itemize}
\end{itemize}

\begin{example}[La Bundesbank face aux chocs pétroliers]
  En 1974, après le premier choc pétrolier, la Bundesbank a resserré sa politique
  monétaire pour lutter contre l'inflation, malgré la récession. Hausse des taux en
  1973-1974~: pour contrer l'inflation post-choc pétrolier, les taux ont été portés à
  des niveaux élevés (par exemple, le taux de prêt marginal a atteint 13~\% en 1974)~;
  assouplissement en 1975-1976~: face à la récession, baisse des taux, mais sans
  sacrifier la lutte contre l'inflation~; nouveau durcissement en 1979, après le
  deuxième choc pétrolier, pour éviter une spirale inflationniste.
\end{example}

\subsubsection{Zoom : effet d'éviction des dépenses privées par les dépenses publiques}

Contexte~: une politique de dépenses publiques financées par emprunt.

\begin{definition}[Effet d'éviction]
  Effet d'éviction = la hausse des dépenses publiques financées par emprunt réduit la
  hausse des investissements privés via la hausse de $r$.
\end{definition}

\begin{remark}[Formulation du cours]
  «~Effet d'éviction = la hausse des investissements publics réduit la hausse des
  investissements privés via la hausse de $r$~».
\end{remark}

Mécanisme~:
\begin{itemize}
  \item hausse de $G$ $\Rightarrow$ hausse de la demande $\Rightarrow$ croissance~;
  \item mais hausse de la demande de monnaie par l'État (demande d'emprunt)
    $\Rightarrow$ hausse de $r$ $\Rightarrow$ réduit un peu les investissements privés
    (et les financements privés) $\Rightarrow$ réduit un peu la production
    $\Rightarrow$ réduit un peu l'effet multiplicateur sur la croissance.
\end{itemize}

\begin{example}[Une relance budgétaire avec accompagnement monétaire]
  Question~: si accompagnement monétaire, l'effet multiplicateur est-il plus fort ou
  plus faible~?\\
  Plus fort. A relance budgétaire with a monetary expansion is the first cell of the
  policy-mix table at the end of the chapter: «~Relance puissante~: hausse PIB
  (multiplicateur max) et $r$ stable (pas d'effet d'éviction)~». The money created
  keeps $r$ from rising, so the effet d'éviction of the borrowing on private
  investment disappears; under changes fixes the same création monétaire is the
  «~effet expansionniste supplémentaire~» of the case of changes fixes above.
\end{example}

\subsubsection{Comment l'inflation a été brisée~?}

\begin{itemize}
  \item Le SME oblige la France à adopter une politique de désinflation pour éviter les
    dévaluations périodiques et pour une croissance saine et durable.
  \item Un changement de politique économique~: «~le grand tournant en 1983~» = passage
    d'une logique de relance keynésienne à une politique de rigueur (plan Delors en
    1983, puis \textbf{désindexation des salaires} (\textit{wage deindexation}) en
    1982-1983)~;
    \begin{itemize}
      \item renforcée par la signature de l'Acte unique en février 1986 (la stabilité
        monétaire est censée stimuler le commerce intra-européen en réduisant les
        incertitudes pour les entreprises souhaitant accroître leurs exportations)~;
      \item soutenue par le \textbf{contre-choc pétrolier} (\textit{oil counter-shock})
        en 1986 (\textbf{désinflation importée} (\textit{imported disinflation})).
    \end{itemize}
  \item Une baisse des taux d'intérêt réels et une crédibilité retrouvée~: la Banque de
    France a gagné en autonomie, et les marchés ont cru en la volonté de l'État de
    lutter contre l'inflation.
\end{itemize}

\subsection{1983-1998 : pourquoi une politique de rigueur monétaire avec des changes fixes requiert un système monétaire unique~?}

After 1983 France chose rigueur and the stability of the franc against the mark. The
SME crisis of 1992--93, which broke out months after the traité de Maastricht had been
signed, showed that changes fixes between currencies with separate central banks stay
exposed to attaques spéculatives however rigorous the policy, and confirmed the move to
the single currency the treaty had planned: the Union économique et monétaire (UEM) and
the BCE, with budget rules attached.

\subsubsection{Politique de franc fort «~ou~» de désinflation compétitive 1983-1998}

\begin{itemize}
  \item Référence au modèle allemand.
  \item Les gouvernements successifs pendant 15 ans privilégient la lutte contre
    l'inflation dans une perspective de long terme~:
    \begin{itemize}
      \item des plans d'austérité (ex.~: plan Delors en 1983, plan Bérégovoy en 1988)
        pour réduire le déficit~: réduction des déficits publics par la hausse des PO et
        le contrôle des dépenses de santé, baisse des dépenses publiques (gel des
        embauches, réduction des subventions)~;
      \item priorité à la stabilité et à la rigueur monétaire plutôt qu'à la relance~:
        stabilité du franc par rapport au mark $\Rightarrow$ \textbf{politique de franc
        fort} (\textit{strong-franc policy}), tenue par tous les temps, y compris au
        début des années 90 lors de la crise du SME, afin de ne pas ruiner la
        crédibilité difficilement acquise~; maintien d'un taux $r$ élevé.
    \end{itemize}
  \item Inflation~: 13,6~\% en 1980 $\Rightarrow$ 5,8~\% en 1985.
\end{itemize}

\begin{definition}[Désinflation compétitive (\textit{competitive disinflation})]
  Désinflation compétitive (aidée par le contre-choc pétrolier en 1986) $\Rightarrow$ si
  monnaie s'apprécie alors baisse du coût des importations $\Rightarrow$ «~désinflation
  importée~» et gains en compétitivité.
\end{definition}

The course gives this name to the whole 1983--1998 policy of the title: hold the franc
against the mark and keep French inflation below that of France's partners, so that
competitiveness is regained without resorting to the periodic dévaluations of the franc
within the SME; the désinflation importée of 1986 helped it.

\subsubsection{Pourquoi une nouvelle crise du SME en 1992-93 malgré la politique de rigueur~?}

\begin{itemize}
  \item Dévaluations répétées~: entre 1979 et 1987, le franc a été dévalué contre le
    mark, ce qui a affaibli la crédibilité du SME.
  \item La réunification allemande (1990) entraîne des transferts budgétaires publics
    massifs $\Rightarrow$ hausse de la demande $\Rightarrow$ inflation $\Rightarrow$
    politique restrictive de la Bundesbank $\Rightarrow$ la hausse des taux d'intérêt
    attire les capitaux vers le mark $\Rightarrow$ appréciation du mark $\Rightarrow$
    affaiblissement des autres monnaies du SME, qui sont obligées de suivre une
    politique monétaire restrictive aussi $\Rightarrow$ récession en Europe en 1993
    (recul du PIB, chômage en hausse).
\end{itemize}

\begin{remark}[Formulation du cours]
  «~$\Rightarrow$ plus grave récession (croissance 1~\% et chômage en hausse) en Europe
  depuis la Seconde Guerre mondiale.~»
\end{remark}

\begin{itemize}
  \item Alimentation des anticipations de dévaluations ou de sortie du SME
    $\Rightarrow$ attaques spéculatives~:
    \begin{itemize}
      \item le week-end du 12-13 septembre 1992, la lire est dévaluée de 7~\%~; le
        16 septembre, la livre est suspendue du mécanisme de change du SME, puis la lire
        le 17.
    \end{itemize}
\end{itemize}

\begin{remark}[Formulation du cours]
  «~Le 13 septembre 1992, la livre est dévaluée de 7,5~\%~; quelques jours plus tard,
  elle sort du SME ainsi que la Lire.~»
\end{remark}

\begin{itemize}
  \item En juillet 1993, attaque spéculative contre le franc $\Rightarrow$ la BdF a
    augmenté ses taux d'intérêt (jusqu'à 13~\% en 1993) pour défendre le franc, avec un
    achat massif de francs contre des devises sur les marchés des changes $\Rightarrow$
    pas de dévaluation, car réforme du SME (= accord de Bruxelles d'août 1993)
    $\Rightarrow$ les marges de fluctuation autour du cours pivot sont élargies à
    $\pm$~15~\%. L'élargissement est tel qu'on ne peut plus véritablement parler de
    fixité des changes entre les monnaies européennes.
\end{itemize}

Bilan~: problème pour maintenir la crédibilité du système $\Rightarrow$ motivation à la
création de l'UEM et d'une banque centrale européenne (BCE), déjà prévues par le
\textbf{traité de Maastricht} (\textit{Maastricht Treaty}) signé en 1992, et à
l'indépendance de la BdF (1994).

\begin{remark}[Formulation du cours]
  «~Bilan~: Problème pour maintenir la crédibilité du système $\Rightarrow$
  motivation à la création de l'UEM et d'une banque centrale européenne (BCE) et à
  l'indépendance de la BdF (1994) $\Rightarrow$ Traité de Maastricht~».
\end{remark}

Changes flexibles were no escape either, as the United Kingdom found in 1985.

\begin{example}[Les limites des changes flexibles : le Royaume-Uni en 1985]
  The case, point by point:
  \begin{itemize}
    \item Les changes flexibles n'empêchent pas les attaques spéculatives et les fortes
      dépréciations. Ce système permet à la monnaie de s'ajuster en fonction de l'offre
      et de la demande, mais expose le pays à des pressions spéculatives et à des
      fluctuations brutales, surtout en période de déséquilibres économiques ou de crise
      de confiance.
    \item Dépréciation de 20~\% de la £ en 1 an.
    \item Origines~: des déséquilibres économiques structurels (déficit commercial
      chronique dans les années 1980, avec une industrie peu compétitive et une
      dépendance aux importations~; la politique monétaire restrictive de la Banque
      d'Angleterre, hausse des taux d'intérêt pour lutter contre l'inflation, a attiré
      des capitaux étrangers, mais a aussi pesé sur la croissance et alourdi la dette
      publique)~; la spéculation contre la livre (les marchés financiers ont anticipé
      une faiblesse prolongée de la livre, notamment en raison des tensions commerciales
      avec les États-Unis~; le dollar était alors très fort).
    \item Conséquences économiques~: inflation importée (la dépréciation a renchéri le
      coût des importations, pétrole et biens manufacturés, alimentant l'inflation)~;
      perte de confiance (les investisseurs étrangers ont réduit leurs placements au
      Royaume-Uni, limitant les financements pour l'économie).
    \item Apprentissage~: la crise de 1985 a montré les limites d'un régime de change
      flexible en cas de déséquilibres économiques profonds. Elle a aussi souligné
      l'importance de la crédibilité des politiques économiques~: la méfiance des
      marchés envers la livre reflétait un manque de confiance dans la capacité du
      gouvernement à maîtriser l'inflation et les déficits.
    \item De plus, les changes flexibles introduisent une incertitude sur les prix des
      biens échangés entre pays européens, ce qui freine les investissements et le
      commerce, mais aussi la signature de contrats à long terme (ex.~: contrats
      industriels ou énergétiques).
  \end{itemize}
\end{example}

\subsubsection{Le traité de Maastricht de 1992}

Signed in 1992, the traité de Maastricht entered into force on 1 November 1993. Its
provisions are now articles 121, 123, 126 and 140 of the traité sur le fonctionnement
de l'Union européenne (TFUE), the current numbering of the provisions that come from
the treaty.

\begin{remark}[Formulation du cours]
  The four articles (121, 123, 126 and 140) are cited as those of the «~Traité
  Maastricht (1992)~».
\end{remark}

\begin{itemize}
  \item La \textbf{coordination des politiques économiques} (\textit{economic policy
    coordination}) (article 121)~: le Conseil fixe, sur recommandation de la Commission,
    les \textbf{grandes orientations de la politique économique} (\textit{broad
    economic policy guidelines}, GOPE) des États membres et de l'Union européenne, surveille
    les politiques économiques menées par les États et leur conformité avec ces
    orientations, et peut adresser des recommandations à un État membre qui ne les
    respecte pas.
  \item Le financement des États membres (article 123)~: l'article 123 interdit aux
    banques centrales, européenne et nationales, d'accorder tout type de crédits aux
    organismes publics des États membres de l'Union et d'acquérir «~directement~»
    auprès d'eux des «~instruments de leur dette~».
  \item Stabilité de l'inflation, des taux de change et des taux d'intérêt de long terme
    (article 140)~:
    \begin{itemize}
      \item la stabilité des prix (taux d'inflation ne dépassant pas de plus de 1,5
        point celui des trois États les plus stables)~;
      \item la \textbf{stabilité du taux de change} (\textit{exchange-rate stability})
        (maintien dans les marges de fluctuation du SME)~;
      \item la \textbf{convergence des taux d'intérêt à long terme}
        (\textit{convergence of long-term interest rates})~: le taux nominal moyen à
        long terme ne dépasse pas de plus de 2 points de pourcentage celui des trois
        États membres, au plus, présentant les meilleurs résultats en matière de
        stabilité des prix.
    \end{itemize}
  \item Les \textbf{déficits excessifs} (\textit{excessive deficits}) (article 126)~:
    «~les États membres évitent les déficits publics excessifs~», et la
    \textbf{discipline budgétaire} (\textit{fiscal discipline}) est respectée si~:
    \begin{itemize}
      \item le rapport entre le déficit public effectif ou prévu et le PIB est inférieur
        à 3~\%, à moins qu'il n'ait diminué de manière substantielle et constante et ne
        se rapproche de cette référence, ou que le dépassement soit limité, exceptionnel
        et temporaire~;
      \item le rapport entre la dette publique et le PIB est inférieur à 60~\%, à moins
        qu'il ne diminue suffisamment et ne s'en rapproche à un rythme satisfaisant.
    \end{itemize}
  \item Si un État membre ne semble pas respecter la discipline budgétaire, la
    Commission examine le solde structurel (et la relation avec les dépenses publiques
    d'investissement) $\Rightarrow$ si déficit excessif, le Conseil adresse des
    recommandations à l'État concerné pour qu'il mette fin à la situation de déficit
    excessif dans un délai donné $\Rightarrow$ en cas de non-respect, «~mise en
    demeure~» $\Rightarrow$ si l'État membre ne se conforme pas à cette mise en demeure,
    le Conseil peut décider qu'il fasse auprès de l'Union «~un dépôt d'un montant
    approprié sans intérêts~» ou lui «~imposer des amendes d'un montant approprié~».
  \item Les décisions et recommandations du Conseil sont prises sans tenir compte du
    vote de l'État membre concerné.
\end{itemize}
These five conditions --- inflation, taux de change, long-term rates, déficit public and
dette publique --- are the critères de Maastricht.

\subsubsection{1993-1997 : politique budgétaire restrictive pour entrer dans la zone euro}

Mesures~:
\begin{itemize}
  \item réduction des dépenses publiques~: gel ou baisse des budgets de certains
    ministères, report de projets coûteux~;
  \item hausse des impôts~: augmentation de la contribution sociale généralisée (CSG)
    en 1996, hausse de la TVA de 18,6~\% à 20,6~\% en 1995~;
  \item réduction des dépenses sociales~: réforme des retraites (allongement de la durée
    de cotisation pour une retraite à taux plein, passée de 37,5 à 40 annuités, ou
    report de l'âge légal) et mesures pour réduire le déficit de la Sécurité sociale
    (indexation des pensions sur les prix plutôt que sur les salaires)~;
  \item privatisations~: vente de 60 entreprises publiques (ex.~: France Télécom, Total,
    Air France) pour réduire la dette (gain = 150 milliards d'euros équivalents
    actuels).
\end{itemize}

\begin{remark}[Privatisations et dette publique]
  Les privatisations n'augmentent pas mécaniquement la dette publique au moment où
  elles ont lieu, mais elles peuvent indirectement contribuer à son augmentation à
  moyen ou long terme, notamment en réduisant les recettes de l'État et en fragilisant
  sa capacité à financer des services publics~: elles transfèrent des revenus stables
  (comme les péages autoroutiers) vers le secteur privé, et les entreprises publiques
  (comme EDF, la SNCF ou les autoroutes) génèrent souvent des bénéfices ou des
  dividendes reversés à l'État, recettes récurrentes que leur privatisation supprime.
  Exceptions~: les privatisations peuvent réduire la dette si les entreprises publiques
  sont chroniquement déficitaires (ex.~: Air France dans les années 1990), ou si le
  produit de la vente est utilisé pour rembourser la dette (ex.~: privatisations en
  Grèce après 2010).
\end{remark}

\begin{itemize}
  \item Accélération de la rigueur des mesures pour atteindre les cibles, et la BdF
    augmente son taux d'intérêt après la crise du franc en 1995 (doute des marchés sur
    la crédibilité de la France).
  \item $\Rightarrow$ Passage à l'euro~: les parités bilatérales sont annoncées en mai
    1998 et les taux de conversion fixés irrévocablement le 31 décembre 1998~; l'euro
    devient la monnaie (scripturale) le 1\textsuperscript{er} janvier 1999~; ses
    billets et pièces remplacent le franc le 1\textsuperscript{er} janvier 2002.
\end{itemize}

\begin{remark}[Formulation du cours]
  «~1998~: Passage à l'euro (fixation irrévocable des parités en mai 1998)~»~;
  elsewhere «~Création de la zone euro en 1999~» and «~Passage à l'euro (2002)~: la
  France abandonne le franc pour adopter la monnaie unique~».
\end{remark}

\subsubsection{Pacte de Stabilité et de Croissance : 1997, réformé en 2024}

The \textbf{Pacte de Stabilité et de Croissance} (\textit{Stability and Growth
Pact}, PSC) dates from 1997. The trajectoire de dépenses, the reduction of the dette and the
sanctions below are rules of its reformed version, in force since 2024; the convergence
of the déficit structurel towards an objectif de moyen terme is a rule of the framework
that preceded the 2024 reform.

\begin{remark}[Formulation du cours]
  These rules are presented under the heading «~Pacte de Stabilité et de Croissance
  (PSC) (1997)~».
\end{remark}

\begin{itemize}
  \item Les 27 États membres de l'UE s'engagent à respecter le PSC, mais les sanctions
    financières ne concernent que les pays de la zone euro (21 membres depuis l'entrée
    de la Bulgarie en janvier 2026).
  \item Le \textbf{volet préventif} (\textit{preventive arm}) vise à éviter les
    dérapages budgétaires. Les États doivent présenter une fois par an une trajectoire
    de dépenses compatible avec la soutenabilité de leurs finances publiques. La
    Commission européenne évalue les plans nationaux, puis le Conseil de l'Union
    européenne les approuve. Un écart important peut ne pas être pris en considération
    s'il résulte d'une circonstance inhabituelle indépendante de la volonté de l'État
    membre, ou en cas de grave récession affectant la zone euro ou l'ensemble de
    l'Union.
  \item Le \textbf{volet correctif} (\textit{corrective arm}) concerne les États qui
    dépassent les seuils de déficit ou de dette $\Rightarrow$ la Commission peut
    recommander l'ouverture d'une \textbf{procédure de déficit excessif}
    (\textit{excessive deficit procedure}, PDE). Cette procédure oblige l'État concerné à
    corriger sa trajectoire budgétaire dans un délai fixé par le Conseil~: fixation
    d'une trajectoire de convergence du déficit structurel vers un \textbf{objectif de
    moyen terme} (\textit{medium-term objective}, OMT), compris entre 0~\% et 1~\% du
    PIB, ainsi que d'un taux de croissance de référence annuel maximal des dépenses
    publiques.
    \begin{itemize}
      \item Possibilité de dépasser ce taux si compensation par la fiscalité~;
        possibilité de relance budgétaire en cas de crise grave.
      \item Les États doivent présenter une «~trajectoire descendante plausible à moyen
        terme~» de la dette~: elle doit baisser en moyenne d'au moins 1 point de PIB
        par an au-delà de 90~\% du PIB, et de 0,5 point entre 60 et 90~\%. Le déficit
        doit être ramené le plus rapidement possible sous les 3~\% du PIB.
    \end{itemize}
\end{itemize}

\begin{remark}[Formulation du cours]
  «~Les États placés sous PDE doivent la réduire de 1~\% par an si elle dépasse 90~\%
  du PIB, de 0,5~\% par an si elle se situe entre 60~\% et 90~\% du PIB.~»
\end{remark}

\begin{itemize}
  \item Les sanctions peuvent atteindre 0,05~\% du PIB tous les 6 mois, mais elles ne
    sont pas automatiques si la dette dépasse 60~\% du PIB, dès lors que la trajectoire
    des dépenses nettes permet de placer la dette sur une tendance descendante
    plausible.
  \item Suspensions possibles. Le Pacte a été suspendu entre 2020 et 2023, avec
    l'activation de la \textbf{clause dérogatoire générale} (\textit{general escape
    clause}) pendant la pandémie de Covid-19 puis dans le contexte de la guerre en
    Ukraine. Les règles sont redevenues applicables en 2024 et permettent des
    investissements (transition écologique, numérique, sécurité énergétique et
    défense). Cette flexibilité passe par une \textbf{clause dérogatoire nationale}
    (\textit{national escape clause}), qui permet à un État de s'écarter temporairement
    de sa trajectoire budgétaire sans être considéré comme en infraction, à condition
    que l'écart soit lié à des dépenses de défense et ne compromette pas la
    soutenabilité de la dette. Elle peut couvrir une période de quatre ans, de 2025 à
    2028, dans la limite de 1,5~\% du PIB par an. 18 clauses nationales ont été
    activées~: quinze États en juillet 2025, puis l'Allemagne en octobre 2025,
    l'Autriche en février 2026 et l'Espagne en juin 2026. La France n'en fait pas
    partie.
\end{itemize}

\subsubsection{Que contient le plan budgétaire français~?}

France's plan under these rules is its \textbf{plan budgétaire et structurel à moyen
terme} (\textit{medium-term fiscal-structural plan}, PSMT).
\begin{itemize}
  \item Le plan vise à ramener le déficit public sous les 3~\% du PIB d'ici 2029. Il
    demande donc un allongement de la période d'ajustement budgétaire de quatre à sept
    ans.
\end{itemize}

\begin{remark}[Formulation du cours]
  «~Il demande donc un allongement de la période d'ajustement budgétaire de quatre à
  cinq ans.~»
\end{remark}

\begin{itemize}
  \item Le déficit public doit descendre progressivement jusqu'à 2,8~\% en 2029, et la
    dette publique culminer vers 116,5~\% en 2027 avant de revenir à 115,8~\% en 2029,
    un niveau toujours bien supérieur au plafond de 60~\% inscrit dans les règles de
    l'UE.
  \item Ce plan repose sur une trajectoire de maîtrise des dépenses et sur plusieurs
    engagements de réformes et d'investissements, notamment dans les domaines du marché
    du travail, de la recherche, de la transition énergétique, de la simplification
    administrative, de l'industrie verte et de l'évaluation de la dépense publique.
  \item Le PSMT français a été validé le 26 novembre 2024 par la Commission européenne,
    une décision que le Conseil de l'Union européenne a entérinée à la majorité
    qualifiée en janvier 2025, lors d'une réunion des ministres européens des Finances.
\end{itemize}

\subsubsection{Les révisions du PSC de 2005 et le TSCG de 2011-12}

Révisions du PSC~:
\begin{itemize}
  \item objectifs = équilibrer discipline budgétaire et croissance en évitant des
    politiques d'austérité trop brutales~;
  \item chaque État doit présenter une trajectoire budgétaire sur plusieurs années,
    adaptée à sa situation~;
  \item indicateur = déficit structurel, puis extension à la croissance des dépenses
    nettes.
\end{itemize}

\begin{definition}[Déficit structurel (\textit{structural deficit})]
  \textbf{Déficit structurel}~: le déficit «~corrigé des variations conjoncturelles~».
\end{definition}

\begin{definition}[Dépenses nettes (\textit{net expenditure})]
  Dépenses nettes = dépenses publiques $-$ charges d'intérêts de la dette $-$ certaines
  mesures temporaires $-$ les dépenses liées aux programmes européens financés par
  l'Union $-$ les dépenses nationales de cofinancement de programmes européens.
\end{definition}

This is the déficit structurel met at the start of the chapter, in the pair
«~Stabilisateurs / déficit structurel~»: the deficit left once the conjoncture is
corrected for.\\

Le \textbf{pacte budgétaire européen} (\textit{Fiscal Compact}), ou \textbf{traité sur
la stabilité, la coordination et la gouvernance} (\textit{Treaty on Stability,
Coordination and Governance}, TSCG) de 2012, impose aux États membres de l'UE de maintenir un
déficit structurel maximal de 0,5~\% du PIB à moyen terme. Certains pays, comme
l'Allemagne, ont intégré cette règle dans leur constitution (via la Schuldenbremse),
avec un objectif de déficit structurel proche de 0~\% (0,35~\% du PIB pour
l'Allemagne).

\subsubsection{Est-ce que la désinflation compétitive est un succès en 1998~?}

Oui~!
\begin{itemize}
  \item Baisse durable et contrôle de l'inflation~: de 13,6~\% en 1980 à 0,7~\% en 1998
    (2~\% en 97).
  \item Croissance et déficit public. The 1997--1998 values of both are in the table
    of the 1997--2001 tournant expansionniste below.
  \item 15 ans de rigueur malgré les alternances politiques.
  \item 15 ans pour construire la crédibilité de la politique de rigueur.
\end{itemize}
Mais\ldots
\begin{itemize}
  \item Chômage élevé (inflation \emph{vs} chômage~!), dont 40~\% de chômeurs de longue
    durée.
  \item Perte de souveraineté monétaire.
  \item Taux d'intérêt élevé $\Rightarrow$ hausse du coût de la dette publique, et
    limite les investissements privés.
  \item Un report des réformes structurelles (retraites, fiscalité) qui auraient pu
    améliorer la compétitivité à long terme.
  \item Déclin de l'emploi industriel~: en 1998, l'emploi industriel en France
    représentait 24~\% de l'emploi total (contre 30~\% en 1980). Le nombre d'emplois
    industriels avait baissé de 1,5 million entre 1980 et 1998. La part de l'industrie
    dans le PIB était passée de 24~\% en 1980 à 18~\% en 1998.
  \item Partage «~inégal~» de la richesse~: entre 1983 et 1998, les salaires réels ont
    progressé moins vite que la productivité.
  \item Pauvreté en hausse~: le nombre de personnes sous le \textbf{seuil de pauvreté}
    (\textit{poverty line}) (60~\% du revenu médian) est passé de 10~\% en 1980 à 13~\%
    en 1998.
  \item Augmentation des inégalités~: les 10~\% les plus riches détenaient 25~\% des
    revenus en 1998, contre 22~\% en 1980.
\end{itemize}

\subsubsection{Nouvelles relations entre les politiques conjoncturelles dans la zone euro : le policy mix}

\begin{itemize}
  \item Avec la création de la BCE --- créée le 1\textsuperscript{er} juin 1998, elle
    conduit la politique monétaire unique depuis le 1\textsuperscript{er} janvier 1999
    --- la politique monétaire est gérée au niveau européen (perte de souveraineté
    nationale).
\end{itemize}

\begin{remark}[Formulation du cours]
  «~Avec la création de la BCE en 1998~»~; elsewhere «~la BCE (créée en 1999)~».
\end{remark}

The BCE, whose seat is in Francfort:
\begin{itemize}
  \item les décisions en matière de politique monétaire pour la zone euro sont prises
    par le \textbf{conseil des gouverneurs de la BCE} (\textit{ECB Governing Council}),
    dans lequel sont représentés les gouverneurs de toutes les banques centrales de
    l'Eurosystème~;
  \item elle est indépendante du pouvoir politique~;
  \item ses objectifs~: définir et mettre en œuvre la politique monétaire de la zone
    euro~; conduire la politique de change de l'euro~; détenir et gérer les réserves de
    change des États membres~; assurer le bon fonctionnement des systèmes de paiement~;
    et, depuis 2014 seulement, avec le mécanisme de surveillance unique (MSU), la
    supervision des grandes banques et l'agrément des établissements de crédit.
\end{itemize}

\begin{remark}[Formulation du cours]
  «~Supervision des principales banques (agrément institution de crédit)~» is listed
  among the «~5 objectifs~» of the BCE at its creation in 1998.
\end{remark}

\begin{definition}[Politique monétaire unique (\textit{single monetary policy})]
  Politique monétaire unique~: la Banque centrale européenne (BCE) fixe les taux
  d'intérêt pour toute la zone euro.
\end{definition}

\begin{itemize}
  \item La BCE suit une politique monétaire conventionnelle (avec les mêmes outils que
    les banques centrales nationales).
  \item Changes flexibles €/autres devises $\Rightarrow$ focus sur le contrôle de
    l'inflation (et non sur le taux de change).
  \item Si peu de risque d'inflation, la politique monétaire peut être expansionniste
    pour favoriser la relance.
  \item La politique budgétaire au niveau national $\Rightarrow$ toujours souveraine,
    mais gestion et contrôle des déficits et des dettes par l'UE pour vérifier le
    respect des critères de Maastricht (le PSC).
\end{itemize}
The policy-mix exercise at the end of the chapter sets the two policies side by side.
What makes a single central bank effective is its crédibilité:
\begin{itemize}
  \item une banque centrale (plus) crédible s'inscrit dans un \textbf{cercle vertueux de
    crédibilité} (\textit{virtuous circle of credibility})~: la crédibilité passée ancre
    les anticipations d'inflation à la cible, réduisant le risque d'une spirale
    inflationniste. Cela réduit l'ampleur des hausses de taux d'intérêt nécessaires
    pour faire baisser l'inflation, ainsi que les coûts associés pour l'activité
    économique~;
  \item la crédibilité permet également de gagner du temps. Lorsqu'une banque centrale
    est imparfaitement crédible, elle doit agir immédiatement pour calmer toute tension
    inflationniste latente et prouver en permanence qu'elle vise à stabiliser
    l'inflation. Mais une banque centrale crédible peut réagir avec un certain décalage
    tout en atteignant un ancrage nominal~;
  \item c'est la réaction moyenne à l'inflation à moyen terme, et non la réaction
    immédiate, qui prime~: une réaction trop rapide et trop forte risque de créer une
    récession et d'alimenter les anticipations d'inflation, et une réaction trop lente
    rend plus difficile de combattre l'inflation. Une réaction moyenne permet donc de
    lisser l'action de la BCE et de l'adapter tout en confortant sa crédibilité.
\end{itemize}

\subsubsection{1997-2001 : pourquoi un tournant expansionniste~?}

\begin{itemize}
  \item La France, sous l'impulsion du gouvernement Jospin (1997-2002), adopte une
    politique keynésienne pour relancer la croissance et réduire le chômage~:
    \begin{itemize}
      \item relance budgétaire~: augmentation des dépenses publiques (éducation, santé,
        infrastructures) et baisse des impôts pour les ménages modestes~;
      \item politique de l'emploi active~: mise en place des 35 heures (loi Aubry en
        1998 et 2000) pour stimuler l'emploi, bien que son impact économique soit
        débattu.
    \end{itemize}
  \item Baisse des taux d'intérêt~: la Banque de France, puis la BCE à partir de 1999,
    réduisent progressivement les taux pour soutenir l'activité économique.
  \item $\Rightarrow$ complémentarité du policy «~mix~» $\Rightarrow$ efficacité
    maximale.
\end{itemize}

The effects, France against Germany, in \%; the taux de chômage are those
standardised by the Organisation de coopération et de développement économiques
(OCDE):
\begin{center}
  \small
  \begin{tabular}{llccccc}
    \toprule
    & & 1997 & 1998 & 1999 & 2000 & 2001 \\
    \midrule
    Taux de croissance du PIB en volume & France & 1,9 & 3,3 & 3,2 & 3,2 & 2,6 \\
    & Allemagne & 1,4 & 2,1 & 1,6 & 3,0 & 2,2 \\
    \addlinespace
    Solde des administrations publiques & France & $-3{,}0$ & $-2{,}7$ & $-1{,}6$ &
      $-1{,}3$ & $-0{,}5$ \\
    rapporté au PIB & Allemagne & $-2{,}7$ & $-2{,}1$ & $-1{,}4$ & $-1{,}5$ & $-1{,}7$
      \\
    \addlinespace
    Taux de chômage au sens du BIT & France & 12,3 & 11,8 & 11,2 & 9,5 & 8,5 \\
    (données standardisées par l'OCDE) & Allemagne & 9,9 & 9,3 & 8,6 & 7,9 & 7,9 \\
    \addlinespace
    Consommation des ménages en volume & France & 0,6 & 2,7 & 2,8 & 2,5 & 2,8 \\
    \bottomrule
  \end{tabular}
\end{center}
Under this policy mix France grows faster than Germany every year from 1997 to 2001,
its déficit public falls from 3,0~\% to 0,5~\% of PIB, its taux de chômage from 12,3~\%
to 8,5~\%, and household consumption in volume, up 0,6~\% in 1997, rises by about
2,7~\% a year from 1998.

\subsection{2007-2026 : pourquoi la succession des crises a augmenté l'endettement en France~?}

From 2007 the shocks follow one another, and each brings public support that raises
the endettement; the relances are followed by spending cuts and tax rises that do not
bring the déficit public down.

\subsubsection{Succession de chocs économiques}

\begin{itemize}
  \item 2007-2008~: crise des subprimes $\Rightarrow$ spéculation immobilière
    $\Rightarrow$ crise financière (défaillance de la régulation bancaire aux États-Unis)
    et crise économique mondiales (sortie de crise mi-2009).
\end{itemize}

\begin{example}[La crise des subprimes, 2007-2008]
  The first shock, step by step:
\begin{itemize}
  \item Contexte~: une politique de taux d'intérêt bas menée de 2001 à 2004 par la
    banque centrale des États-Unis (la FED) pour contrer la récession de 2001 due à
    l'explosion de la bulle Internet.
  \item $\Rightarrow$ Les banques américaines (et des sociétés financières non bancaires
    américaines faiblement régulées) octroient de nombreux crédits immobiliers, souvent
    à taux variables, aux ménages à faibles revenus. Ces crédits sont appelés
    «~subprimes~», car ils présentaient un risque élevé de non-remboursement. Ces
    crédits immobiliers étaient garantis par le bien immobilier acheté (\textbf{crédit
    hypothécaire} (\textit{mortgage loan})). Les ménages étaient incités à parier sur
    une revente à terme de leurs actifs immobiliers pour rembourser les crédits
    souscrits. Une fois les crédits octroyés, les banques se séparaient facilement de
    ces crédits en les transformant en titres, vendus à des investisseurs sur les
    marchés financiers~: c'est la \textbf{titrisation} (\textit{securitisation}).
  \item Cependant, la FED, confrontée aux tensions liées à la forte croissance (hausse
    des prix, risque de bulles spéculatives), décide d'augmenter son taux d'intérêt
    directeur, qui passe de 1 à 5~\% entre 2004 et 2006. Cette augmentation se
    répercute sur les taux d'intérêt variables des crédits subprimes. Les mensualités
    de crédit progressent fortement, au point que les ménages se trouvent dans
    l'impossibilité d'y faire face. Plusieurs millions de ménages américains se
    retrouvent en défaut de paiement et doivent renoncer à leur bien immobilier.
  \item La mise en vente des nombreux logements saisis par les banques et sociétés
    financières contribue à la chute des prix de l'immobilier. Les investisseurs
    financiers détenteurs des titres associés aux crédits subprimes connaissent
    d'importantes pertes. La crise immobilière devient une crise financière. Les
    agences de notation dégradent les notes qu'elles accordent aux acteurs financiers
    et aux obligations adossées aux crédits hypothécaires subprimes. Certains
    établissements, notamment Lehman Brothers, font faillite.
  \item Le \textbf{risque systémique} (\textit{systemic risk}) se matérialise par le jeu
    des interdépendances entre acteurs, secteurs et économies~; un effet de contagion
    menace la stabilité du système financier mondial. La méfiance s'installe entre les
    banques, paralysant le marché interbancaire. Elles refusent de se prêter de
    l'argent et la moindre rumeur entraîne des flux de retraits significatifs de la
    part de nombreux investisseurs. Fragilisées, elles accordent aussi moins de crédits
    aux ménages et entreprises.
  \item La crise financière se propage à l'économie. La consommation et l'investissement
    diminuent, provoquant une chute du PIB, une explosion du chômage et une baisse des
    prix, donc un risque de déflation. En outre, l'état des finances publiques se
    dégrade et la dette publique augmente fortement.
\end{itemize}
\end{example}

The shocks that follow:
\begin{itemize}
  \item 2008-2012~: \textbf{crise des dettes souveraines} (\textit{sovereign debt
    crisis}) en Europe (Grèce, Italie, Espagne, France) $\Rightarrow$ déficit et
    endettement public $\Rightarrow$ risque de crise financière européenne (sortie de
    crise en 2013).
  \item 2015-2020~: risque de déflation $\Rightarrow$ \textbf{politique monétaire non
    conventionnelle} (\textit{unconventional monetary policy}), l'assouplissement
    quantitatif (QE).
  \item 2020~: crise Covid $\Rightarrow$ choc d'offre $\Rightarrow$ économie au ralenti
    et soutiens aux entreprises $\Rightarrow$ rupture d'approvisionnement et inflation
    post-Covid.
  \item 2022~: guerre en Ukraine $\Rightarrow$ pénurie agricole et énergétique
    $\Rightarrow$ inflation générale.
  \item 2026~: guerre en Iran $\Rightarrow$ rupture d'approvisionnement en
    pétrole/gazole/urée $\Rightarrow$ inflation énergétique (essence) et risque sur
    l'alimentation et les biens manufacturés (hausse du coût du diesel $\Rightarrow$
    hausse du fret maritime $\Rightarrow$ hausse des produits importés d'Asie). The
    barrel of West Texas Intermediate (WTI), around 60~\$ until February 2026, passes
    100~\$ from March to May 2026.
  \item Bilan~: envol des soutiens publics («~quoi qu'il en coûte~», bouclier tarifaire
    énergétique, bouclier gazole, etc.)~; montée des incertitudes mondiales
    $\Rightarrow$ réduire les dépendances $\Rightarrow$ recherche de souveraineté pour
    être moins impacté par les chocs externes.
\end{itemize}

\subsubsection{La crise des dettes souveraines en zone euro, 2008-2012}

\textbf{Origines et pays concernés.}
\begin{itemize}
  \item Ralentissement de la \textbf{croissance structurelle} (\textit{structural
    growth}), c'est-à-dire à long terme $\Rightarrow$ moins de recettes fiscales et, en
    même temps, hausse des dépenses publiques $\Rightarrow$ déficits publics.
  \item Excès de liquidités dans le monde (politique monétaire accommodante)
    $\Rightarrow$ investissements/spéculation accrus.
  \item Création de la zone euro en 1999~: entrée de capitaux importants durant une
    décennie dans les pays dits périphériques, alimentant un endettement excessif tant
    public que privé.
  \item Crise de 2008 $\Rightarrow$ hausse des dépenses publiques pour soutenir
    l'activité $\Rightarrow$ hausse des déficits publics.
  \item La crise a d'abord concerné la dette publique grecque avant de se propager à
    l'Irlande (2010), au Portugal (2011), à Chypre, à l'Italie puis à l'Espagne (tous
    les trois en 2012).
\end{itemize}

\textbf{Pourquoi des attaques spéculatives~?}
\begin{itemize}
  \item Dégradation des notes sur les dettes $\Rightarrow$ hausse des taux (renforcée
    par la faillite de certaines banques nationales) $\Rightarrow$ effet boule de
    neige.
    \begin{itemize}
      \item Il faut gérer les déficits et pas seulement le niveau d'endettement
        (déficit grec/PIB = 15,4~\% en 2009~; dettes grecques/PIB = 126~\% en 2009).
      \item Concernant l'Italie, c'est l'insuffisance de son plan d'austérité pour
        réduire sa dette (126~\% du PIB).
    \end{itemize}
\end{itemize}

\textbf{Gestion et conséquence pour la politique monétaire en Europe.}
\begin{itemize}
  \item Défaut d'un pays $\Rightarrow$ hausse des taux d'intérêt de long terme
    $\Rightarrow$ krach obligataire~!
  \item En mai 2010, un prêt de 110 milliards d'euros est accordé à la Grèce par la zone
    euro et le FMI $\Rightarrow$ politique d'assainissement des comptes publics
    (réduction des transferts sociaux, récession).
  \item Le \textbf{Mécanisme européen de stabilité} (\textit{European Stability
    Mechanism}, MES), créé en 2012, prend le relais du Fonds européen de stabilité
    financière (FESF) et du Mécanisme européen de stabilisation financière (MESF), créé
    en 2010. Il est permanent et dispose d'une capacité de prêt
    initiale de 500 milliards d'euros, alimentée par les contributions financières de
    chaque pays.
\end{itemize}

\begin{remark}[Formulation du cours]
  «~Le Mécanisme européen de stabilisation financière (MESF) en 2012 remplace le Fonds
  européen de stabilité financière (FESF). Il est permanent et dispose d'une capacité
  de prêt initiale de 500 milliards d'euros.~»
\end{remark}

\subsubsection{Le risque de déflation, 2015-2020}

\begin{itemize}
  \item Origines~: pendant la crise de 2008, deux canaux de transmission sont bloqués~:
    \begin{itemize}
      \item le \textbf{canal des taux d'intérêt} (\textit{interest-rate channel})
        $\Rightarrow$ fortes baisses des taux d'intérêt entre 2007 et 2009 sur le marché
        monétaire, mais aversion au risque $\Rightarrow$ arrêt des prêts
        interbancaires~;
      \item le \textbf{canal des crédits} (\textit{credit channel}) $\Rightarrow$ les
        banques commerciales ont conservé les liquidités injectées par les banques
        centrales sous forme de réserves plutôt que d'accorder des prêts
        supplémentaires.
    \end{itemize}
  \item Entre 2015 et 2020, l'inflation dans la zone euro est restée durablement sous la
    cible des 2~\% (en France, moyenne de 0,5~\%), car~:
    \begin{itemize}
      \item chute des prix de l'énergie~: entre 2014 et 2016, le prix du pétrole a chuté
        de plus de 50~\% (passant de 99~\$ à 44~\$ le baril) $\Rightarrow$ réduction de
        l'inflation~;
      \item faible croissance des salaires~: malgré une légère amélioration du marché du
        travail, les hausses de salaires sont restées limitées, limitant la pression
        inflationniste par la demande~;
      \item globalisation et concurrence accrue~: la mondialisation a maintenu les prix
        des biens manufacturés à un niveau bas, en raison de la concurrence des pays à
        bas coûts (Chine, Europe de l'Est, etc.)~;
      \item les ménages ont préféré épargner plutôt que consommer, par crainte de
        l'avenir (chômage, retraite\ldots)~;
      \item autres facteurs~: Brexit (2016), guerres commerciales États-Unis/Chine à
        partir de 2018.
    \end{itemize}
  \item Politique monétaire ultra-accommodante~: taux négatifs + \textbf{assouplissement
    quantitatif} (\textit{quantitative easing}, QE)~: à partir de mars 2015, rachat
    massif d'actifs par la BCE (obligations d'État et privées) pour injecter des
    liquidités dans l'économie, mais impact limité (c'est le Covid et la guerre en
    Ukraine qui ont relancé l'inflation~!).
\end{itemize}
The BCE's taux directeur stayed at about 0~\% from 2015 until mid-2022. The BCE ran a
«~politique ZIRP (Zero Interest Rate Policy) pour faire baisser les taux de long terme
en maintenant «~longtemps~» des taux courts nuls/très faibles entre 2015 et 2021 en
Europe~».

\subsubsection{La politique monétaire non conventionnelle et le Quantitative Easing}

\textbf{Quantitative Easing (QE).}
\begin{itemize}
  \item Montant massif et programmé~: la BCE annonce des programmes d'achats massifs.
  \item Durée longue et prédéfinie~: les achats s'étalent sur plusieurs années.
  \item Effet systémique~: les achats massifs font baisser les rendements des
    obligations (ex.~: les OAT françaises), ce qui réduit le coût de la dette pour les
    États et stimule l'investissement.
  \item Ciblage explicite~: la BCE annonce des programmes dédiés~:
    \begin{itemize}
      \item programme APP (Asset Purchase Programme)~: de 2015 à 2022, la BCE a acheté
        pour 2\,600 milliards d'euros d'actifs, dont des OAT~;
      \item programme PEPP (Pandemic Emergency Purchase Programme)~: en 2020-2022, la
        BCE a acheté 1\,850 milliards d'euros d'actifs, dont des OAT, pour faire face à
        la pandémie~;
      \item selon une étude de l'Observatoire français des conjonctures économiques
        (OFCE), les banques centrales de la zone euro ont acheté 58~\% des emprunts
        publics émis en 2020 (73~\% pour la France) et détenaient 25~\% du stock de
        dettes publiques à la fin de l'année (23~\% pour la BdF)~;
      \item 488 milliards d'euros (au coût amorti) détenus par la BdF au 30 juin 2026
        pour les deux programmes.
    \end{itemize}
\end{itemize}

\textbf{Cadre juridique et légitimité.}
\begin{itemize}
  \item Achats d'OAT (hors QE)~: ces achats s'inscrivent dans le mandat traditionnel de
    la BCE (stabilité des prix, soutien à la politique économique de l'UE) $\to$ base
    légale~: TFUE (art. 127).
  \item Quantitative Easing~: le QE est une interprétation extensive du mandat de la
    BCE, justifiée par des circonstances exceptionnelles (crise financière, pandémie)
    $\to$ base légale~: la jurisprudence de la Cour de justice de l'UE (CJUE). La CJUE
    a validé le QE (programme PSPP, Public Sector Purchase Programme) par l'arrêt Weiss du 11 décembre 2018~; l'arrêt
    Gauweiler (2015) portait sur le programme OMT (Outright Monetary Transactions).
\end{itemize}

\begin{remark}[Formulation du cours]
  «~Base légale~: Jurisprudence de la Cour de justice de l'UE (ex.~: arrêt Gauweiler en
  2015 validant le QE)~».
\end{remark}

\begin{notation}
  Here OMT is the BCE's Outright Monetary Transactions programme, not the objectif de
  moyen terme (OMT) of the PSC.
\end{notation}

\textbf{Risques du QE.}
\begin{itemize}
  \item Inflation.
  \item Hausse de l'endettement des États.
  \item Risque de bulle financière et de menace sur la stabilité (hausse de la prise de
    risques des investisseurs, car il est facile de s'endetter pour «~spéculer~») =
    effets de bord~; + risque de crise de la dette en 2022 en Europe (hausse des taux
    avec la sortie du QE).
\end{itemize}

How QE lowers the cost of borrowing for a fragile State goes through the prime de
risque.

\begin{definition}[Prime de risque (\textit{risk premium})]
  La prime de risque (ou spread) correspond à la différence de rendement entre la dette
  d'un État et un actif sans risque (comme les obligations allemandes ou américaines).
\end{definition}

Le QE agit sur cette prime via plusieurs canaux~:
\begin{itemize}
  \item baisse des taux longs~: en achetant massivement des obligations d'État (ou
    d'autres actifs), la banque centrale réduit l'offre de titres disponibles sur le
    marché. Cela fait monter leur prix et donc baisser leur rendement (taux d'intérêt).
    Une baisse des taux longs réduit mécaniquement la prime de risque, car les
    investisseurs exigent moins de compensation pour détenir cette dette~;
  \item \textbf{effet de signal} (\textit{signalling effect}), ou forward guidance~: le
    QE envoie un signal fort aux marchés sur la volonté de la banque centrale de
    maintenir des conditions financières accommodantes. Cela réduit l'incertitude et la
    perception de risque de défaut, surtout pour les États perçus comme fragiles~;
  \item liquidité accrue~: en injectant des liquidités dans le système bancaire, le QE
    facilite le refinancement des États et réduit les tensions sur les marchés
    obligataires. Cela peut atténuer les pics de prime de risque en période de crise~;
  \item \textbf{effet richesse} (\textit{wealth effect}) et demande~: en stimulant
    l'activité économique (via la baisse des taux et l'augmentation du crédit), le QE
    peut améliorer les perspectives de croissance et de recettes fiscales des États.
    Cela réduit indirectement le risque perçu de défaut.
\end{itemize}

\subsubsection{Différentes mesures de relance et de soutien, 2007-2025}

Des mesures à court terme pour endiguer les crises et éviter une récession prolongée,
stabiliser et soutenir l'activité économique\ldots
\begin{itemize}
  \item aides pendant le Covid~;
  \item plan de sauvetage bancaire (Crédit Agricole, BNP Paribas)~: 360 milliards
    d'euros en 2008~; plus le plan Sarkozy (26 milliards d'euros)~;
  \item allègement de l'ISF (hausse du seuil à 1,3 million d'euros, plafonnement à
    75~\% du revenu), suppression de l'ISF (remplacé par l'IFI en 2018) et flat tax
    (30~\% sur les revenus du capital)~;
  \item baisse de l'IS (de 33,3~\% en 2017 à 25~\% en 2022), et suppression du crédit
    d'impôt pour la compétitivité et l'emploi (CICE) et de la taxe professionnelle~;
  \item augmentation du SMIC (+1,5~\% en 2019)~;
  \item baisse puis suppression de la taxe d'habitation sur la résidence principale~;
  \item mesures discrétionnaires pour protéger le pouvoir d'achat et l'activité
    économique~: 81 milliards d'euros en 2022-23 (boucliers tarifaires, aides aux
    entreprises énergivores).
\end{itemize}
\ldots et des mesures à moyen et long terme~:
\begin{itemize}
  \item plan France Relance (2021-22)~: 100 milliards d'euros de plans de relance
    post-Covid, pour la transition écologique (30~Md€), la compétitivité des
    entreprises (34~Md€) et la cohésion sociale/pouvoir d'achat (36~Md€)~;
  \item Grand emprunt national (2010)~: 35 milliards d'euros~; plan d'investissement de
    57 milliards d'euros (2018-2022)~; plan France 2030 (2021)~: 54 milliards d'euros
    sur 5 ans, pour financer des projets d'avenir (recherche, numérique, transition
    écologique, modernisation des infrastructures, souveraineté).
\end{itemize}

\subsubsection{Après les relances, des réductions des dépenses\ldots}

\begin{itemize}
  \item Gel des dépenses publiques et suppression de 150\,000 postes de fonctionnaires
    (2010), gel des salaires de la fonction publique (2011-12).
  \item Réforme des retraites (2010)~: report de l'âge légal de départ de 60 à 62 ans
    et de l'âge du taux plein de 65 à 67 ans. Cette réforme, très contestée, est
    adoptée sans référendum~; la réforme des retraites par points est ensuite suspendue
    en 2020~; puis réforme des retraites (2023)~: passage à 64 ans.
\end{itemize}

\begin{remark}[Formulation du cours]
  «~Cette réforme, très contestée, est adoptée sans référendum mais échec de celle de
  2015 puis Réforme des retraites (2023)~: passage à 64 ans.~»
\end{remark}

\begin{itemize}
  \item Réforme de l'assurance-chômage (2019)~: durcissement des conditions
    d'indemnisation.
  \item Hausse des impôts~:
    \begin{itemize}
      \item sur les revenus (2011, 2012, 2014) et modification des impôts sur les hauts
        revenus (plus de 150\,000~€ par an)~;
      \item contributions exceptionnelles sur les hauts revenus (2010, 2012, 2013, 2020,
        2021, 2023, 2024, 2025, 2026) et sur les bénéfices de plus de 1 milliard d'euros
        (2022-23), pour des gains de 10 à 20 milliards d'euros par an au maximum~;
      \item réduction des niches fiscales~;
      \item taxation du diesel (hausse progressive de la TICPE), hausse du malus
        écologique~;
      \item hausse de la taxe sur les transactions financières (de 0,3~\% à 0,4~\%)
        (2025)~;
      \item hausse du taux normal de TVA de 19,6~\% à 20~\% (2014)~; the 21,2~\% rate
        voted in 2012 never came into force.
    \end{itemize}
\end{itemize}

\begin{remark}[Formulation du cours]
  «~Hausse de la TVA (passée de 19,6~\% à 21,2~\%) (2010-12)~»
\end{remark}

\begin{itemize}
  \item Accroissement des contrôles pour fraude fiscale.
\end{itemize}

Bilan~: pas de réduction des déficits publics, et hausse de la dette publique. In the
first quarter of 2026 (Eurostat), la France, l'Italie et la Grèce ont des ratios entre
117,6 et 143,5~\% du PIB, far above the 60~\% reference.

\subsubsection{Bilan 2026 : nombreux pays sous procédure de déficits excessifs}

\begin{itemize}
  \item Au 15 juillet 2026, dix États membres sont sous procédure de déficit excessif~:
    Autriche, Belgique, Bulgarie, Finlande, France, Hongrie, Italie, Pologne, Roumanie
    et Slovaquie.
  \item La France est concernée depuis 2024 (et l'a été de 2009 à 2018). Le Conseil a
    recommandé qu'elle mette fin à sa situation de déficit excessif d'ici 2029, avec une
    trajectoire limitant la progression nominale de ses dépenses nettes sur la période
    2025-2029.
  \item Selon les données Eurostat (avril 2026)~:
    \begin{itemize}
      \item le déficit public s'élevait en 2025 à 2,9~\% du PIB dans la zone euro et à
        3,1~\% du PIB dans l'Union européenne~;
      \item la dette publique atteignait 87,8~\% du PIB dans la zone euro et 81,7~\% dans
        l'Union européenne~;
      \item la France faisait partie des pays les plus endettés de l'Union, avec une
        dette publique de 115,6~\% du PIB à la fin de 2025, derrière la Grèce et
        l'Italie.
    \end{itemize}
\end{itemize}

\subsection{Conclusion}

\begin{itemize}
  \item Rôle crucial de la crédibilité des mesures de politiques économiques
    (monétaires et budgétaires).
  \item La zone euro réduit les risques de spéculation, mais la hausse des dettes
    souveraines peut les réveiller~: selon Eurostat (T1 2026), Grèce (143,5~\%), Italie
    (138,9~\%), France (117,6~\%, mais plus grand montant total de l'UE), Belgique
    (109,1~\%) et Espagne (101,6~\%).
  \item La situation actuelle résulte de plusieurs crises et chocs $\Rightarrow$
    nombreuses variables interdépendantes $\Rightarrow$ les prévisions sont compliquées
    $\Rightarrow$ «~peu d'accalmie économique~»~!
  \item Dès 2024, la Cour des comptes estime que la trajectoire n'est pas compatible
    avec une stabilisation rapide de la dette sans ajustement additionnel sur la
    période 2025-2027.
  \item En 2026, en France~: croissance molle, endettement non contrôlé.
\end{itemize}
The question this leaves, «~Comment (peut-on) sauver la France~?~», is the subject of
the next chapter.

\subsection{Questions de compréhension : expliquer les mécanismes}

\begin{example}[Inflation et taux de change]
  Question~: quelle relation entre le taux d'inflation et le taux de change~?\\
  Une demande globale soutenue par les dépenses publiques peut pousser les prix à la
  hausse. Si l'inflation en France est plus élevée que chez ses partenaires commerciaux
  (comme l'Allemagne), les biens français deviennent moins compétitifs, ce qui réduit
  les exportations et affaiblit la demande pour le franc $\Rightarrow$ baisse de la
  valeur du franc. Bilan~: inflation différentielle sur les autres pays $\Rightarrow$
  dépréciation. The fuite des capitaux of the différentiel d'inflation above is the
  second channel to the same result.
\end{example}

\begin{example}[Déficit extérieur et réserves de change]
  Question~: quelle relation entre le déficit extérieur et la pression sur les réserves
  de change~?\\
  Les dépenses publiques élevées, en stimulant la demande intérieure, peuvent aussi
  aggraver le déficit commercial~: si les importations augmentent (pour répondre à la
  demande intérieure) plus vite que les exportations, le solde extérieur se dégrade.
  Conséquence~: la France doit puiser dans ses réserves de change pour acheter des
  francs et soutenir sa monnaie, ce qui est coûteux et insoutenable à long terme.
  Effet~: une fois les réserves épuisées, la Banque de France ne peut plus défendre le
  franc, ce qui conduit à une dévaluation forcée.
\end{example}

\subsection{Quel «~policy mix~» idéal pour quelle situation~?}

\begin{definition}[Policy mix]
  Le policy mix, c'est la combinaison de la politique budgétaire (État~: dépenses
  publiques $G$, impôts $T$) et de la politique monétaire (banque centrale~: masse
  monétaire $M \to$ taux d'intérêt $r$). Les deux politiques ne produisent pas les
  mêmes effets, ne passent pas par les mêmes canaux, et surtout peuvent se renforcer
  ou s'annuler.
\end{definition}

\begin{example}[Les quatre configurations du policy mix]
  Question~: quel «~policy mix~» idéal pour quelle situation~?\\
  Each combination of the two policies has its use:
  \begin{center}
    \small
    \begin{tabularx}{\linewidth}{>{\raggedright\arraybackslash}p{3.2cm}
        >{\raggedright\arraybackslash}X >{\raggedright\arraybackslash}X}
      \toprule
      & Expansion monétaire & Restriction monétaire \\
      \midrule
      Expansion budgétaire & Relance puissante~: hausse du PIB (multiplicateur max) et
        $r$ stable (pas d'effet d'éviction). & Freiner l'inflation en soutenant la
        croissance, mais risque d'annuler la croissance. \\
      \addlinespace
      Restriction budgétaire & Réduire le déficit sans réduire la croissance. & Juguler
        une inflation très forte~; risque élevé de récession. \\
      \bottomrule
    \end{tabularx}
  \end{center}
\end{example}

\begin{remark}[Questions de révision]
  Three questions are left without an answer.
  \begin{itemize}
    \item En quoi un déficit commercial peut-il entraîner une hausse de la dette
      publique~?
    \item Signification de «~Le taux d'intérêt de la BCE augmente de 25 points~»~?
    \item Trouver un exemple ou cas d'usage pour chacune des 4 configurations possibles
      du policy mix.
  \end{itemize}
\end{remark}

\begin{remark}[Concepts clés]
  The concepts to know from this chapter:
  \begin{itemize}
    \item Les perfectionnements successifs du système monétaire et financier en Europe
      pour appréhender les chocs~: systèmes monétaires (Serpent, SME, UEM) et contrôle
      des déficits publics (PSC).
    \item Politiques monétaires conventionnelle et non conventionnelle.
    \item Policy mix et coordination en Europe~: la combinaison optimale des deux
      politiques.
    \item Politique contra-cyclique vs.\ pro-cyclique et impacts sur les déficits publics
      et la dette publique.
    \item Relations inflation, balance commerciale, taux d'intérêt, taux de change,
      dévaluation, attaque spéculative.
    \item Le triangle d'incompatibilité de Mundell-Fleming.
  \end{itemize}
\end{remark}
